\chapter{Физическое происхождение фермионной петли для логарифма определителя KMS}
\label{ch:v9-kms-physical-fermion-loop-origin}

\section{Проверяемая физическая альтернатива}

После запретительный результат Stückelberg-механизма остаётся проверить, может ли требуемый
barrier возникнуть как обычный однопетлевой determinant уже существующих
физических fermions. Target равен
\begin{equation}
 B_{\rm target}=-\log\det R_\theta-\log\det R_\kappa,
 \qquad
 Z_{\rm target}=\det R_\theta\det R_\kappa.
 \label{eq:v9-kms-physical-loop-target}
\end{equation}
Здесь
\begin{equation}
 R_\theta=\operatorname{diag}(r_{\theta s},r_{\theta a},
 r_{\theta t},r_{\theta t},r_{\theta t}),
 \quad
 R_\kappa=\operatorname{diag}(r_{\kappa s},r_{\kappa a},
 r_{\kappa t},r_{\kappa t},r_{\kappa t}).
 \label{eq:v9-kms-physical-loop-type-operators}
\end{equation}

\section{Физическая multiplicity creation-cell}

Creation-cell имеет одну configuration-source line и одну пятимерную
target-мультиплету:
\begin{equation}
 \mathcal H_{\rm cr}=\mathbb C\Omega\oplus V_{\rm type},
 \qquad \dim_{\mathbb C}V_{\rm type}=1+1+3=5.
 \label{eq:v9-kms-physical-loop-carrier}
\end{equation}
Соответствующий projector имеет
\begin{equation}
 P_{\rm type}=\operatorname{diag}(0,1,1,1,1,1),
 \qquad \rank P_{\rm type}=5.
 \label{eq:v9-kms-physical-loop-projector}
\end{equation}
Даже если интерпретировать все пять target-lines как одну complex fermion
multiplet, её quadratic kernel даёт один determinant
\begin{equation}
 \int D\bar\psi D\psi\,
 e^{-\bar\psi K\psi}=\det K.
 \label{eq:v9-kms-physical-loop-gaussian}
\end{equation}

\section{Degree obstruction для block-linear loop}

Каждый type determinant имеет weighted degree five:
\begin{equation}
 \det R_\theta=r_{\theta s}r_{\theta a}r_{\theta t}^{3},
 \qquad
 \det R_\kappa=r_{\kappa s}r_{\kappa a}r_{\kappa t}^{3}.
 \label{eq:v9-kms-physical-loop-single-determinants}
\end{equation}
Для любого homogeneous block-linear kernel на одной пятимерной multiplet
общая rescaling даёт
\begin{equation}
 \det\!\left[\lambda(aR_\theta+bR_\kappa)\right]
 =\lambda^5\det(aR_\theta+bR_\kappa).
 \label{eq:v9-kms-physical-loop-linear-degree}
\end{equation}
Target одновременно масштабируется как
\begin{equation}
 \det(\lambda R_\theta)\det(\lambda R_\kappa)
 =\lambda^{10}\det R_\theta\det R_\kappa.
 \label{eq:v9-kms-physical-loop-target-degree}
\end{equation}
Поэтому одна обычная linear fermion multiplet имеет determinant capacity
$1/2$: она не воспроизводит два независимых blocks.

\section{Почему Real doubling не является второй species}

Антисимметричный Real lift одного kernel имеет вид
\begin{equation}
 \mathcal A_J(R)=
 \begin{pmatrix}0&R\\-R&0\end{pmatrix},
 \qquad
 \det\mathcal A_J(R)=(\det R)^2.
 \label{eq:v9-kms-physical-loop-real-lift}
\end{equation}
Physical Pfaffian half-count возвращает один determinant,
\begin{equation}
 \operatorname{Pf}\mathcal A_J(R)=\pm\det R,
 \label{eq:v9-kms-physical-loop-pfaffian}
\end{equation}
а не independent $\det R_\theta\det R_\kappa$. Попытка назначить двум
Real-партнёрам разные blocks нарушает exchange compatibility. Для exact
witness
\begin{equation}
 R_\theta^{(w)}=\operatorname{diag}(1,2,3,3,3),
 \quad
 R_\kappa^{(w)}=\operatorname{diag}(2,1,2,2,2)
 \label{eq:v9-kms-physical-loop-real-witness}
\end{equation}
коммутатор block-оператора с Real exchange имеет
\begin{equation}
 \rank[J_{\rm ex},R_\theta^{(w)}\oplus R_\kappa^{(w)}]=10.
 \label{eq:v9-kms-physical-loop-real-mismatch}
\end{equation}

\section{Два algebraic routes и их origin}

Две independent physical multiplets условно дали бы
\begin{equation}
 K_{\rm double}=R_\theta\oplus R_\kappa,
 \qquad \dim K_{\rm double}=10,
 \qquad \det K_{\rm double}=Z_{\rm target}.
 \label{eq:v9-kms-physical-loop-double}
\end{equation}
Но в creation-cell зарегистрирована только одна target-мультиплета; вторая
пятимерная physical species является новым носителем datum.

Формально можно остаться на пяти измерениях и записать
\begin{equation}
 K_{\rm comp}=R_\theta R_\kappa,
 \qquad \det K_{\rm comp}=Z_{\rm target}.
 \label{eq:v9-kms-physical-loop-composite}
\end{equation}
Однако это не block-linear inherited kernel. Уже первая компонента содержит
явный mixed coupling
\begin{equation}
 \frac{\partial^2(K_{\rm comp})_{ss}}
 {\partial r_{\theta s}\partial r_{\kappa s}}=1.
 \label{eq:v9-kms-physical-loop-mixed-derivative}
\end{equation}
Ни Hamiltonian creation-frame, ни KMS relation не выводят этот product.
Запись $K_{\rm comp}$ поэтому кодирует требуемый determinant в bare
bilinear и является содержащий целевую подстановку.

\section{Conductance obstruction и Keldysh boundary}

Gap-shape принадлежит Hamiltonian sector, тогда как conductance-shape
задаёт dissipative Kossakowski rates:
\begin{equation}
 H_{\rm cr}=\Delta_sP_s+\Delta_aP_a+\Delta_tP_t,
 \qquad
 \gamma_\alpha^-=\kappa_\alpha,
 \quad \gamma_\alpha^+=\kappa_\alpha e^{-\beta\Delta_\alpha}.
 \label{eq:v9-kms-physical-loop-sector-split}
\end{equation}
В текущем microscopic action нет fermionic bilinear с eigenvalues
$r_{\kappa\alpha}$. Чтобы превратить rates в loop kernel, необходимо
предъявить environment fermions, их spectrum и system--bath coupling.

Schwinger--Keldysh doubling также не является второй particle species:
две contour-ветви описывают одну density evolution, а normalization на
совпадающих sources требует
\begin{equation}
 Z_{\rm SK}[J,J]=1.
 \label{eq:v9-kms-physical-loop-sk-normalization}
\end{equation}
Поэтому branch duplication нельзя считать свободным вторым determinant.

\section{ProofDSL и итог}

LCF certificate проверяет ranks, determinant degrees, Real mismatch,
doubled block и composite-kernel identity:
\begin{equation}
 N_{\rm ProofDSL\ obligations}=11/11,
 \qquad N_{\rm registered\ gates}=56,
 \qquad N_{\rm registered\ obligations}=443.
 \label{eq:v9-kms-physical-loop-proofdsl}
\end{equation}
Итоговые ledgers:
\begin{equation}
 \begin{aligned}
 N_{\rm physical\ target\ carrier}&=5/5,\\
 N_{\rm independent\ determinant\ capacity}&=1/2,\\
 N_{\rm candidate\ parent\ origin}&=0/6,\\
 N_{\rm physical\ logdet\ parent}&=0/1,\\
 N_{\rm physical\ four\ slot\ parent}&=0/1.
 \end{aligned}
 \label{eq:v9-kms-physical-loop-ledger}
\end{equation}

\begin{theorem}[Точная determinant-capacity physical target]
Одна пятимерная physical target-мультиплета несёт один degree-five
fermion determinant. Две независимые copies или один явно multiplicative
kernel алгебраически воспроизводят target degree ten.
\end{theorem}

\begin{nogocriterion}[Physical loop не выводит conductance determinant]
Real/Pfaffian doubling не создаёт independent species, Keldysh doubling
нормирован как один closed-time process, а conductances не входят в
унаследованный Hamiltonian bilinear. Две algebraic реализации требуют либо
новой пятимерной fermion multiplet, либо содержащий целевую подстановку product kernel.
Следовательно, полный physical логарифм определителя parent не выведен.
\end{nogocriterion}

\begin{protocol}\begin{enumerate}
\item physical target multiplet: \textbf{$5$};
\item one-loop linear degree: \textbf{$5$};
\item target degree: \textbf{$10$};
\item determinant capacity: \textbf{$1/2$};
\item Real exchange mismatch witness: \textbf{rank $10$};
\item conditional double: \textbf{$R_\theta\oplus R_\kappa$};
\item conditional composite: \textbf{$R_\theta R_\kappa$};
\item candidate origin: \textbf{$0/6$};
\item ProofDSL: \textbf{$11/11$, registry $56/443$};
\item следующий гейт: \textbf{minimal fermion-bath architecture}.
\end{enumerate}\end{protocol}

Машинный сертификат сохранён в
\path{s2t_v9_endpoint_creation_kms_logdet_physical_fermion_loop_parent_origin_gate_results.json}.

\begin{thebibliography}{9}
\bibitem{v9-kms-physical-loop-uolea}
A.~Angelescu, P.~Huang,
\emph{Integrating Out New Fermions at One Loop},
JHEP \textbf{01} (2021) 049, arXiv:2006.16532.
\bibitem{v9-kms-physical-loop-keldysh}
A.~Kamenev, A.~Levchenko,
\emph{Keldysh technique and non-linear sigma-model: basic principles and
applications}, Adv. Phys. \textbf{58} (2009) 197--319,
arXiv:0901.3586.
\bibitem{v9-kms-physical-loop-sk-brst}
F.~M.~Haehl, R.~Loganayagam, M.~Rangamani,
\emph{Schwinger--Keldysh formalism I: BRST symmetries and superspace},
JHEP \textbf{06} (2017) 069, arXiv:1610.01940.
\end{thebibliography}